% Part: lambda-calculus
% Chapter: introduction
% Section: church-rosser

\documentclass[../../../include/open-logic-section]{subfiles}

\begin{document}

\olfileid{lam}{int}{cr}
\olsection{د چرچ--روسر خاصيت}

\begin{thm}
\ollabel{thm:church-rosser}
فرض کړئ $M$، $N_1$ او $N_2$ داسې ترمونه دي چې $M \red N_1$ او
$M \red N_2$۔ بيا داسې ترم $P$ شته چې $N_1 \red P$ او $N_2 \red P$۔
\end{thm}

\begin{cor}
که $M$ عادي بڼې ته راکمېدے شي، دغه عادي بڼه يوازينۍ ده۔
\end{cor}

\begin{proof}
که $M \red N_1$ او $M \red N_2$ وي، د مخکينۍ قضيې له مخې داسې
ترم~$P$ شته چې $N_1$ او $N_2$ دواړه~$P$ ته راکميږي۔ که $N_1$ او
$N_2$ دواړه په عادي بڼه کښې وي، دا يوازې هغه وخت شوني دي چې
$N_1 \ident P \ident N_2$ وي۔
\end{proof}

په پاى کښې، دوه ترمونه $M$ او~$N$ \emph{$\beta$-معادل}، يا يوازې
\emph{معادل}، هغه وخت بولو چې دواړه يوه ګډ ترم ته راکميږي؛ په بله
وينا، داسې $P$ شته چې $M \red P$ او $N \red P$۔ دا په
$M \equal[\beta] N$ ليکو۔ د \olref{thm:church-rosser} په کارولو
څېړلے شئ چې $\equal[\beta]$ د معادلتوب اړيکه ده، او دا اضافي
خاصيت هم لري چې د هر $M$ او~$N$ دپاره، که $M \red N$ يا
$N \red M$ وي، نو $M \equal[\beta] N$۔ (په حقيقت کښې، ښودل کېدے
شي چې $\equal[\beta]$ د دې خاصيت لرونکې \emph{تر ټولو کوچنۍ}
د معادلتوب اړيکه ده۔)

\end{document}
